\section{Conclusions}
\label{sec:conclusions}

This study evaluated an executable cross-shard commit protocol through an
integrated, reproducible evidence workflow under one frozen experimental
protocol. Bounded property verification, mutation-based validation,
implementation-model trace conformance, verification-cost characterization,
and artifact reproduction provide complementary evidence rather than
cumulative proof, and no individual technique is claimed as novel in
isolation.

For RQ1, 350 of 420 measured bounded executions completed, while 70 TLC-large
executions were censored by timeout, with no property violations observed
among completed configurations. This is bounded evidence rather than an
unbounded correctness proof. For RQ2, all 10 predefined scientific mutants
were detected, giving a mutation score of 1.0 for the evaluated catalogue.
The result demonstrates sensitivity to those selected defect classes, not
completeness against arbitrary faults.

For RQ3, 600 trace executions produced the expected bounded conformance
outcomes: valid traces were accepted and deliberately corrupted traces were
rejected with the expected diagnostics. This does not establish refinement or
behavioral equivalence. For RQ4, completed Alloy profiles showed increasing
elapsed time and memory use, while all measured TLC-large executions reached
timeout. These observations characterize the evaluated verification envelope,
not a nonlinear or asymptotic scaling law or absolute TLC-versus-Alloy
superiority.

Taken together, these layers form a reproducible argument in which bounded
verification, targeted defect sensitivity, trace conformance, and cost
characterization address distinct limitations. Their combination strengthens
the evaluation without extending any layer beyond its stated evidential scope.

The analytical artifact was reproduced from a separate clone and workspace
under a separate operating-system user on the same native Linux host. The
attempt completed 10/10 validation gates and matched 32/32 expected
analysis-artifact hashes. This establishes process separation and deterministic
reconstruction under the documented environment, not hardware independence, a
complete rerun of all 1,272 scheduled tasks, or scientific correctness from
hash agreement alone.

Further work can extend the evaluated bounds, broaden the scientific mutant
catalogue and trace scenarios, and repeat reproduction on an independent
physical machine. These extensions would strengthen external validity while
preserving the bounded interpretation of the present evidence.
